\section{Conclusion}
\label{sec:conclusion}

This work hides a payload in a threaded discussion by choosing a reply target and a comment topic, then asking an LLM to write an ordinary comment.
Capacity is whatever those two lists uniquely decode.
The design does not insert bits into next-token sampling, and it does not need white-box logits.

An independent LLM-judge run on the 244-post unique-ZLG subset~\cite{lin2024zero} found ZLG~\cite{lin2024zero} harder to pick in a same-thread lineup (28.3\% versus 51.6\% standout detection) and our method higher on register-fit (3.193 versus 2.791, paired sign-test $p=0.00374$).
The source campaign, run with one retry rather than the default six, succeeded on 2{,}543 of 4{,}044 attempts, all by audit-assisted reconstruction from sender-side metadata.
A historical 2026-07-30 configuration, averaged over 100 posts, reported an available fixed-width selection budget of 15.14 bits/comment for our method versus 16.00 reported bits for ZLG~\cite{lin2024zero}, and GPT-2 perplexity of 68.15 versus 53.70 (ours versus ZLG~\cite{lin2024zero}), with 554 of 627 project-side attempts succeeding. These bit figures do not measure independently recovered project-side payload; the ZLG~\cite{lin2024zero} reveal used sender token IDs and externally supplied payload length.
Those two artifacts answer different questions.
Neither is a blind, capacity-matched ranking, and neither establishes resistance to an adaptive steganalyst.

An exploratory reply-chain probe removed a cross-thread topic mismatch on one frozen post and passed strict decoding on its first frame.
Planning the 64-bit payload required 19 frames, and only 13 of the first 25 posts fit within 20 frames.
This probe does not demonstrate complete payload recovery and shows the capacity cost of restricting topics to the selected reply chain.
Future work should test complete recovery with the normal receiver, freeze a symmetric manifest, score GPT-2 and divergence on that freeze, decode without sender audit files, and test dedicated steganalysis and editing attacks.
